\subsection{伪反射}

定义 1. 向量空间 \(s\) 的端同态 \(V\) 称为伪反射，如果 \(1 - s\) 的秩为 1。

设 \(s\) 是 \(V\) 中的一个伪反射，\(D\) 是 \(1 - s\) 的像。按定义，\(D\) 的维数为 1；于是给定 \(a \neq 0\) 中的 \(D\)，存在 \(a^*\) 上的一个非零线性形式 \(V\)，使得 \(x - s(x) = \langle x, a^* \rangle\) . \(a\)，对于所有 \(x \in V\)。

反之，给定 V 中的 \(a \neq 0\) 和 V 上的非零线性形式 \(a^{*} \neq 0\)，公式

\[
s _ {a, a ^ {*}} (x) = x - \langle x, a ^ {*} \rangle . a \quad (x \in V)\tag{1}
\]

定义一个伪反射 \(s_{a,a^*}\)；\(1 - s_{a,a^*}\) 的像由 \(a\) 生成，而 \(1 - s_{a,a^*}\) 的核是 \(V\) 的超平面，它由那些 \(x\) 满足 \(\langle x, a^* \rangle = 0\) 的点组成。如果 \(V^*\) 是 \(V\) 的对偶空间，立即可见，\(s_{a^*,a}\) 是 \(s_{a,a^*}\) 的转置，是 \(V^*\) 的伪反射，由公式

\[
s _ {a ^ {*}, a} (x ^ {*}) = x ^ {*} - \langle x ^ {*}, a \rangle . a ^ {*} \quad (x ^ {*} \in V ^ {*}).\tag{2}
\]

如果 \(a\) 是一个非零向量，带向量 \(a\) 的伪反射是任意伪反射 \(s\)，使 \(a\) 属于 \(1 - s\) 的像。伪反射 \(s\) 的超平面是 \(1 - s\) 的核，即由满足 \(x\) 的向量 \(s(x) = x\) 组成的集合。

命题 1. 设 G 是一个群，\(\rho\) 是 G 在向量空间 V 上的一个不可约线性表示；假设 G 中存在元素 g，使得 \(\rho(g)\) 是一个伪反射。

\begin{enumerate}
\def\labelenumi{(\roman{enumi})}
\item
与 \(\rho(\mathbf{G})\) 对易的 V 的每个端同态都是一个仿射伸缩，且 \(\rho\) 绝对不可约。
\item
  假设 V 是有限维的。设 B 是 V 上在 \(\rho(G)\) 下不变的非零双线性型。则 B 非退化，或者对称或者斜对称，并且 V 上在 \(\rho(G)\) 下不变的每个双线性型都与 B 成比例。
\end{enumerate}

设 \(u\) 是与 \(V\) 对易的 \(\rho(G)\) 的一个端同态。设 \(g\) 是 \(G\) 中的一个元素，使得 \(\rho(g)\) 是伪反射，令 \(D\) 为 \(1 - \rho(g)\) 的像。由于 \(D\) 的维数为 1 且 \(u(D) \subset D\)，存在 \(\alpha\) 属于 \(K\)，使得 \(u - \alpha.1\) 在 \(D\) 上为零；于是 \(N\) 的核 \(u - \alpha.1\) 是 \(V\) 的一个在 \(\rho(G)\) 下不变的向量子空间，并且由于包含 \(D\) 而非零；由于 \(\rho\) 不可约，\(N = V\) 且 \(u = \alpha.1\)。(i) 的第二部分由第一部分以及《代数》，第 VIII 章，§13，第 4 号，命题 5 的推论得出。

设 N（分别地，\(N'\)）是 V 中由满足对所有 y 属于 V，都有 \(B(x,y)=0\)（分别地，\(B(y,x)=0\)）的 x 组成的子空间；由于 B 在 \(\rho(G)\) 下不变，V 的子空间 N 和 \(N'\) 在 \(\rho(G)\) 下稳定，并且由于 \(B\neq0\) 而不等于 V。由于 \(\rho\) 不可约，\(N=N'=0\)，B 非退化。

由于 \(V\) 是有限维的，\(V\) 上的每个双线性型都可由公式

\[
\mathrm{B} ^ {\prime} (x, y) = \mathrm{B} (u (x), y)\tag{3}
\]

如果 \(B'\) 在 \(\rho(G)\) 下不变，则端同态 u 与 \(\rho(G)\) 对易。事实上，设 x、y 属于 V，g 属于 G；由于 B 和 \(B'\) 都在 \(\rho(G)\) 下不变，有

\[
\begin{array}{r l} & {\mathrm{B} (u (\rho (g) (x)), y) = \mathrm{B} ^ {\prime} (\rho (g) (x), y) = \mathrm{B} ^ {\prime} (x, \rho (g ^ {- 1}) (y))} \\ & {\qquad = \mathrm{B} (u (x), \rho (g ^ {- 1}) (y)) = \mathrm{B} (\rho (g) (u (x)), y),} \end{array}
\]

因此，由于 \(u(\rho(g)(x)) = \rho(g)(u(x))\) 非退化，\(B\)。由 (i)，存在 \(\alpha\) 属于 \(K\)，使 \(u = \alpha.1\)，所以 \(B' = \alpha.B\)。

特别地，我们可以将此应用于双线性型 \(\mathrm{B}^{\prime}(x,y)=\mathrm{B}(y,x)\)；于是对于 V 中所有 x、y，有 \(\mathrm{B}(y,x)=\alpha.\mathrm{B}(x,y)=\alpha^{2}.\mathrm{B}(y,x)\)，而由于 B 非零，有 \(\alpha^{2}=1\)，因此 \(\alpha=1\) 或 \(\alpha=-1\)。所以 B 或者对称，或者斜对称。

